\section{Introduction}\label{sec:intro}

A clinical risk model is judged on two properties. Discrimination is whether it ranks patients
correctly, so that those who have the event were given higher risks than those who do not.
Calibration is whether the risk it reports is the risk the patient carries. A model can rank well and
still report the wrong number, and in medicine the number is what is used: a patient is treated,
screened or discharged when the predicted risk crosses a threshold. A model that predicts $8\%$ for
patients whose true risk is $5\%$ sends the wrong patients over that threshold, however well it
discriminates: discrimination asks only whether the model separates the patients at higher risk from
those at lower risk, not whether the risks it reports are right \citep{vancalster2016hierarchy}. Testing calibration is therefore routine in the
development and validation of clinical prediction models.

For logistic regression with continuous covariates the classical tests do not work. Almost every
patient has a covariate pattern of their own, so the Pearson and deviance statistics are sums of
cells with one observation each, and their chi-squared reference fails. Two repairs followed. One
standardises the Pearson statistic by its moments \citep{mccullagh1985,osius1992normal}. The other, by
\citet{hosmer1980goodness}, sorts the patients by predicted risk, cuts them into ten groups and
compares observed with expected events in each group. Grouping makes the cells large enough for a
reference distribution to exist, and the Hosmer--Lemeshow test is still the default check of
calibration.

Grouping has a cost: a data set of a million patients is reduced to ten numbers before the test sees
it, and the statistic spends one degree of freedom on each group, so its power is spread over every
direction in which a model can fail. Adding groups does not recover the lost information, because
each new group adds a degree of freedom faster than it adds signal, and the verdict can change from
rejection to acceptance with the number of groups alone \citep{paul2013power}. Most variants of the
test repair one of these problems and leave the others. \citet{paul2013power} and \citet{Lai2018}
standardise its power in large samples, \citet{nattino2020assessing} test whether the misfit exceeds a
size that does not grow with $n$, \citet{surjanovic2024generalized} derive a correct reference
distribution for a generalized form of the statistic, and BAGofT \citep{BAGofT2019} learns the
partition from the data instead of fixing it. \citet{hosmer2002goodness} came closest to a directed
grouped test: they weight each record's residual within the groups to gain power against a specified
alternative. None of them lets the partition be refined freely.

Other tests recover power by leaving the grouping behind. \citet{stukel1988generalized} embeds the
logistic link in a two-parameter family and tests the two extra parameters with a score test; in the
comparison of \citet{hosmer1997comparison} it had the best power against a misspecified link and was
among the tests they recommended. The GiViTI calibration belt \citep{nattino2014new,nattino2015new}
fits a polynomial to the calibration curve, and the Liu projection test \citep{liu2024comprehensive}
projects the individual residuals and refers them to a bootstrap. These tests use every record's
residual at its own weight, and that is why they are powerful. It also makes them sensitive to a few
wrong records. Clinical data contain such records: a covariate entered in the wrong units, a value
with a misplaced decimal point, a count typed as a larger but plausible count, which passes any range
check. With a model that is correct for every other patient, one record in a thousand with a
corrupted covariate raises the false-alarm rate of Stukel's score test from $5\%$ to $12\%$, and ten such
records raise it to $63\%$; in the hospital cohort of Section~\ref{sec:realdata}, about twelve miscoded counts among
$35{,}000$ patients raise it to $45\%$. A test that rejects for this reason is reporting a data-entry
error, and the analyst cannot tell it from a miscalibrated model; the two call for opposite actions,
cleaning the data or rebuilding the model.

Two tests keep a form of pooling without a fixed partition. The smoothed residual test of
\citet{le1991goodness} pools residuals over neighbourhoods in covariate space, and BAGofT pools them
over a learned partition. Le Cessie's test is protected against a few wrong records, and BAGofT showed
no rise in a small study (Section~\ref{sec:contamination}), but both pay in computation. Le Cessie's test builds an
$n\times n$ kernel, which on the $35{,}000$ patients of each half of the cohort of
Section~\ref{sec:realdata} would need about $10$ gigabytes, so it was not run there. BAGofT refits the
model more than ten thousand times and takes $392$ seconds for one data set of a thousand patients.

We keep the grouping and project the grouped residuals onto a small
basis of smooth calibration shapes, chosen to include the link departures that Stukel's test targets;
the weights act on groups, never on a single record. The statistic then spends a fixed number of
degrees of freedom however many groups there are, so the partition can be refined as the sample grows
without diluting the test, and its reference distribution needs no resampling or refitting: the test
takes $4$ milliseconds at $n=1000$ and half a second at $n=100{,}000$. We call the test EDGE (Efficient Directed Grouped Examination), the name
it has in the software. Its contributions are:

\begin{enumerate}
\item \textbf{A mechanism.} Pooling on equal-frequency groups bounds the influence that one corrupted
record's covariate can have on a goodness-of-fit statistic, because the record's residual is divided by
the variance of its clean groupmates (Proposition~\ref{prop:bounded}); the bound is exact in external
validation, where a frozen model is checked on new patients. A test loses this protection when it divides
a record's contribution by something that the record or the refit can make small: its own fitted
variance, or the information that the test's direction keeps after the fit. Stukel's test, the cubic
calibration test, an ungrouped test along EDGE's own directions and EDGE's score form all do, and break;
so do grouped tests whose extreme groups are small or that group in covariate space, and a robust fit or
a robust version of Stukel's test restores little
(Sections~\ref{sec:mechanism} and~\ref{sec:contamination}).
\item \textbf{A grouped test whose partition is a setting.} Because its degrees of freedom do not
depend on the number of groups (Theorem~\ref{thm:growingG}), the partition can be refined for power on
clean data and kept coarse for protection when records may be wrong
(Section~\ref{sec:partitionchoice}). In external validation it needs only outcomes and predicted
probabilities, so it applies to any model that reports them, logistic or not; there, at ten groups, its
average power is that of Cox's recalibration test and it keeps its false-alarm rate below $0.10$ where
that test does not. On clean data the test
is comparable in power to Stukel's score
test and the cubic calibration test; unlike them, it keeps its level when a few records are wrong, at
a modest cost in power.
\item \textbf{Evidence, including where the test loses.} A simulation study over six departure families
with every comparison and threshold fixed before it was run, twelve comparator tests including two
resampling tests and four other grouped tests, a study of external validation against the usual tests
there, and a cohort of $70{,}000$ patients in which
EDGE rejects at every partition while the Hosmer--Lemeshow verdict changes with the number of
groups.
\end{enumerate}

Section~\ref{sec:test} defines the test and its null distribution, and Section~\ref{sec:mechanism}
shows why pooling protects it. Section~\ref{sec:design} describes the simulation study,
Section~\ref{sec:power} reports power on clean data and Section~\ref{sec:contamination} under
corrupted records. Section~\ref{sec:realdata} applies the test to the hospital cohort, and
Section~\ref{sec:discuss} says when it is the right test and when another is better.
